\documentclass[11pt,a4paper]{article}
\usepackage[margin=28mm]{geometry}
\usepackage[T1]{fontenc}
\usepackage{lmodern,microtype,amsmath,amssymb,amsthm,mathtools}
\usepackage{enumitem}
\usepackage[hidelinks]{hyperref}
\usepackage[nameinlink,noabbrev]{cleveref}
\setlength{\parindent}{1.2em}
\setlength{\parskip}{2pt}
\setlist[enumerate]{label=\textup{(\roman*)},leftmargin=2em}
\newtheorem{theorem}{Theorem}[section]
\newtheorem{lemma}[theorem]{Lemma}
\newtheorem{proposition}[theorem]{Proposition}
\newtheorem{corollary}[theorem]{Corollary}
\theoremstyle{remark}
\newtheorem{remark}[theorem]{Remark}
\newcommand{\A}{\mathcal A}
\newcommand{\U}{\mathcal U}
\newcommand{\T}{\mathcal T}
\newcommand{\hinv}{h_{\mathrm{inv}}}
\newcommand{\hout}{h_{\mathrm{out}}}
\newcommand{\htrack}{h_{\mathrm{inv}}^+}
\newcommand{\rinv}{r_{\mathrm{inv}}}
\newcommand{\h}{h_{\mathrm{top}}}
\title{What the feasible lift does not determine}
\author{}
\date{P6 research note, 2 October 2026}
\begin{document}
\maketitle

The cone construction in P6 separates the information needed for viability
from the complexity of feasible state trajectories. Here we keep the input
alphabet fixed and compare entire feasible lifts. Their dynamics, input
factors, and state-trajectory factors can be conjugate while all three
control entropy rates vary. A scalar example then tests the common-input
chaos question on an actual control set.

\section{Definitions}

Let $X$ be a metric space, $U$ a compact input space, and
$f:X\times U\to X$ continuous. Inputs are arbitrary sequences
$u\in\U=U^{\mathbb N_0}$. Write $\varphi(k,x,u)$ for the forward solution.
A compact set $Q$ is controlled invariant when each $x\in Q$ has an input
that keeps its trajectory in $Q$ for all nonnegative times.
For compact $K\subset Q$, $\rinv(n,K,Q)$ is the least number of inputs
needed so that each $x\in K$ stays in $Q$ at times $0,\ldots,n$ under one
of those inputs. Set
\[
 \begin{aligned}
 \hinv(K,Q)&=\limsup_{n\to\infty}n^{-1}\log\rinv(n,K,Q),\\
 \hout(K,Q)&=\lim_{\varepsilon\downarrow0}\limsup_{n\to\infty}
       n^{-1}\log\rinv(n,K,Q^\varepsilon),
 \end{aligned}
\]
where $Q^\varepsilon=\{x:d(x,Q)<\varepsilon\}$ uses the ambient metric.
All logarithms are natural.

A tracking catalogue is a finite family $(a_i,u^i)$ whose reference
trajectories remain in $Q$ forever, such that each $x\in K$ satisfies
\[
 \max_{0\leq k\leq n}d(\varphi(k,x,u^i),\varphi(k,a_i,u^i))<\varepsilon
\]
for one member of the family. Its least size is
$\rinv^+(n,\varepsilon,K,Q)$, and $\htrack(K,Q)$ takes the same ordered
limits as $\hout$, using this count. In particular $\hout\leq\htrack$.
Tracking does not require the actual trajectory to be feasible.

The entire feasible lift and its input factor are
\[
 \Lambda_Q=\{(u,x):\varphi(k,x,u)\in Q\text{ for all }k\geq0\},\qquad
 F(u,x)=(\sigma u,f(x,u_0)),\qquad \pi(u,x)=u.
\]
Let $e:\Lambda_Q\to\T_Q$ send a lift point to its entire state trajectory,
and let $S$ be the left shift on $\T_Q\subset Q^{\mathbb N_0}$.
These are compact systems and $eF=Se$.
The uniform input-fiber entropy is
\[
 h_{\rm fib}(F\mid\pi)=\lim_{\varepsilon\downarrow0}\limsup_{n\to\infty}
 \frac1n\log\sup_{u\in\pi(\Lambda_Q)}s_n(F,\varepsilon,\pi^{-1}(u)),
\]
where $s_n$ is the separated-set count for the restricted Bowen metric.

\section{One alphabet, conjugate lifts, and varying rates}

Fix an integer $m\geq2$ and the labelled input alphabet
$\A_m=\{0,\ldots,m-1\}$. Let $1<q\leq m$, and put
\[
 v_j(q)=q-1-\frac{2(q-1)j}{m-1}\qquad(j\in\A_m).
\]
Thus the labels are fixed, although their controlled maps depend on $q$.
Let $R:\mathbb R\to\mathbb R$ be an increasing smooth diffeomorphism with
$R(0)=0$ and $0<R(s)<s$ for $0<s\leq1$. Write $\rho=R'(0)$ and define
\[
 \begin{gathered}
 t(s)=\begin{cases}R(s)/s,&s\ne0,\\ \rho,&s=0,\end{cases}\qquad
 f_j(s,z)=\bigl(R(s),t(s)(qz+sv_j(q))\bigr),\\
 Q=\{(s,z):0\leq s\leq1,\ |z|\leq s\}.
 \end{gathered}
\]
The function $t$ is smooth and positive, since
$t(s)=\int_0^1R'(\theta s)\,d\theta$.

\begin{theorem}\label{thm:realization}
For every $m\geq2$ and every ordered pair
\[
 0<a\leq\log m,\qquad 0\leq b\leq a,
\]
there is a system in this family for which each controlled map is a global
smooth diffeomorphism, $Q$ is controlled invariant, and every compact
$K\subset Q$ of positive area satisfies
\[
 \hinv(K,Q)=a,\qquad \hout(K,Q)=\htrack(K,Q)=b.
\]
The exact logarithmic counts have an actual time limit. The outer and
tracking logarithmic counts have an actual time limit at each fixed
$\varepsilon>0$.

For fixed $m$, all the resulting feasible lifts are conjugate over the
same full input shift. Their state-trajectory factors are also conjugate,
and these conjugacies commute with the trajectory maps $e$. In all cases
\[
 \h(F)=\log m,\qquad h_{\rm fib}(F\mid\pi)=0,\qquad \h(S)=0.
\]
Viewed as measures on the common ambient space
$\A_m^{\mathbb N_0}\times Q$, their invariant probability measures are
exactly $\nu\times\delta_{(0,0)}$, where $\nu$ is shift invariant.
Every such measure has zero conditional Kolmogorov--Sinai entropy over
the input factor.
\end{theorem}

The following interval observation gives the exact upper bound without
an interior viability margin.

\begin{lemma}\label{lem:intervals}
If a finite collection of closed intervals of length $\ell>0$, all
contained in a closed interval of length $L$, covers that interval, a
subcollection of size at most
$2L/\ell+2$ still covers it.
\end{lemma}
\begin{proof}
Take a subcover minimal by inclusion, delete duplicate intervals, and
order the remaining left endpoints $a_1<\cdots<a_N$.
Consecutive intervals meet; otherwise there would be a gap in the
covered interval. If $a_{i+2}\leq a_i+\ell$, the union of intervals $i$
and $i+2$ contains interval $i+1$, contradicting minimality. Thus
$a_{i+2}>a_i+\ell$ whenever these indices occur.
The left endpoints lie in a range of length at most $L-\ell$.
In each parity, successive left endpoints are more than $\ell$ apart,
so each parity has at most $L/\ell+1$ members. Adding the two bounds
proves the assertion.
\end{proof}

\begin{lemma}\label{lem:radial}
The iterates $R^n$ converge uniformly to zero on $[0,1]$ and form an
equicontinuous family there. For each $0<\eta<\rho$ there is $C_\eta$ such
that
\[
 R^n(s)/s\leq C_\eta(\rho+\eta)^n\quad(0<s\leq1).
\]
For every $s_0>0$ there is $c_{\eta,s_0}>0$ such that
\[
 R^n(s)/s\geq c_{\eta,s_0}(\rho-\eta)^n\quad(s_0\leq s\leq1).
\]
\end{lemma}
\begin{proof}
The decreasing sequence $R^n(1)$ has a fixed-point limit. The assumptions
force that limit to be zero, and monotonicity bounds all other iterates
by it. Near zero, $R(s)/s$ lies between $\rho-\eta$ and $\rho+\eta$.
All trajectories enter this neighborhood after one uniform finite time.
Apply these bounds to
$R^n(s)/s=\prod_{k<n}R(R^k(s))/R^k(s)$ and absorb the finitely many earlier
products into constants. Their upper bound is uniform down to $s=0$;
their lower bound is positive on $[s_0,1]$.
For equicontinuity, sufficiently late iterates have uniformly small range,
and the finitely many earlier iterates are uniformly continuous.
\end{proof}

\begin{proof}[Proof of the entropy formulas in \cref{thm:realization}]
The inverse of $f_j$ is given by
\[
 s=R^{-1}(s'),\qquad z=q^{-1}\bigl(z'/t(s)-sv_j(q)\bigr).
\]
It is smooth and globally defined. At positive radius write $z=sy$.
Then $s_k=R^k(s)$ and the normalized state obeys
$y_{k+1}=qy_k+v_{u_k}(q)$.

Let $B=[-1,1]$ and $\psi_j(y)=(y-v_j(q))/q$.
Each interval $\psi_j(B)$ is contained in $B$. The two extreme intervals
reach the endpoints of $B$, and consecutive left endpoints differ by
$2(q-1)/(q(m-1))\leq2/q$, their common length. Hence these intervals
cover $B$, proving viability.
For a word $w$ of length $n$, its viable normalized initial set is exactly
$\psi_{w_0}\cdots\psi_{w_{n-1}}(B)$. These intervals cover $B$ and have
length $2q^{-n}$. Length gives the lower bound $q^n$ for a subcover;
\cref{lem:intervals} gives the upper bound $2q^n+2$.
The same word is viable on every radial slice with the corresponding
normalized initial interval, and the slice $s=1$ gives the converse.
Thus
\begin{equation}\label{eq:exact-count}
 q^n\leq\rinv(n,Q,Q)\leq2q^n+2.
\end{equation}
For a single word the viable vertical section at radius $s$ has length
$2sq^{-n}$, so its total area in $Q$ is $q^{-n}$.
Covering any positive-area $K$ requires at least $|K|q^n$ words, while
\eqref{eq:exact-count} supplies an upper bound. Consequently
$\lim_n n^{-1}\log\rinv(n,K,Q)=\log q$.

Fix $\varepsilon>0$ and $0<\eta<\rho$, and put
$\gamma=q(\rho+\eta)$. A grid in $B$ of at most
$C_\delta\max\{1,\gamma^n\}$ centers ensures
\[
 \max_{0\leq k\leq n}\gamma^k|y-c|<\delta.
\]
Choose an infinite viable normalized input from each center $c$.
Under that input the trajectories from $(s,sy)$ and $(s,sc)$ differ by
$R^k(s)q^k(y-c)$ vertically, whose absolute value is at most
$C_\eta\delta$. Choosing $\delta<\varepsilon/(2C_\eta)$ produces an
outer catalogue with upper rate at most $\max\{0,\log\gamma\}$.

For tracking, equicontinuity in \cref{lem:radial} supplies a finite radial
net $T\subset[0,1]$, independent of $n$, so that each $s$ has $r\in T$
with $\sup_k|R^k(s)-R^k(r)|<\varepsilon/(2\sqrt2)$.
Use references $(r,rc)$ with the chosen viable inputs. If $c_k\in B$
is the reference normalized state, the trajectory difference is
\[
 \bigl(s_k-r_k,(s_k-r_k)c_k+s_kq^k(y-c)\bigr).
\]
Its norm is less than $\varepsilon$, after decreasing $\delta$ if needed.
The number of references is the grid size times $|T|$, giving the same
upper rate. Letting $\eta\downarrow0$ proves both upper bounds
$\max\{0,\log(\rho q)\}$.

For the outer lower bound choose $s_0>0$ so that
$K_0=K\cap\{s\geq s_0\}$ has positive area. At time $n$, a fixed input
maps the initial vertical coordinate affinely with slope
$q^nR^n(s)/s$. Outer-feasible final coordinates lie in the bounded vertical
projection of $Q^\varepsilon$. By \cref{lem:radial}, a single input can
therefore cover at most
$C_{\varepsilon,\eta,s_0}[q(\rho-\eta)]^{-n}$ of the area of $K_0$.
Fubini's theorem and the covering inequality give lower rate at least
$\log(q(\rho-\eta))$. Nonnegativity and $\eta\downarrow0$ give
$\max\{0,\log(\rho q)\}$. The inequality $\hout\leq\htrack$ gives the
tracking lower bound as well. These arguments hold for each fixed
$\varepsilon$, so the two time limits exist at that scale.

Finally take $q=e^a$. If $b<a$, choose $R(s)=e^{b-a}s$; if $b=a$, choose
$R(s)=\operatorname{arsinh}s$. Both are global smooth diffeomorphisms with
the required radial decrease, and the formulas give exactly $(a,b)$.
\end{proof}

\begin{proof}[Proof of the lift and factor assertions]
Let $\Sigma_m=\A_m^{\mathbb N_0}$ and set
\[
 Y_q(u)=-\sum_{k=0}^{\infty}\frac{v_{u_k}(q)}{q^{k+1}}.
\]
The series is uniformly convergent and $|Y_q(u)|\leq1$.
Also $qY_q(u)+v_{u_0}(q)=Y_q(\sigma u)$.
It follows that
\[
 H_{q,R}(u,s)=(u,(s,sY_q(u)))
\]
maps $\Sigma_m\times[0,1]$ into the feasible lift and intertwines $F$ with
$\sigma\times R$. Conversely, a feasible positive-radius trajectory has
bounded normalized states, and
\[
 y=-\sum_{k=0}^{n-1}\frac{v_{u_k}(q)}{q^{k+1}}+q^{-n}y_n
       \longrightarrow Y_q(u).
\]
At zero radius the only state is the apex. Thus $H_{q,R}$ is onto the
entire lift; its input and radial coordinates make it injective.
Compactness gives a homeomorphism.

Any two radial maps $R_1,R_2$ in this construction are conjugate on
$[0,1]$. To see this, choose an increasing homeomorphism
$h_0:[R_1(1),1]\to[R_2(1),1]$ and define
\[
 h(R_1^nt)=R_2^nh_0(t)\quad
   (R_1(1)\leq t\leq1,\ n\geq0),\qquad h(0)=0.
\]
Adjacent definitions agree at endpoints. The shells cover $(0,1]$, and
their images tend to zero, proving continuity at zero.
The resulting increasing homeomorphism satisfies $hR_1=R_2h$.
Hence
$H_{q_2,R_2}(\mathrm{id}\times h)H_{q_1,R_1}^{-1}$ is an input-preserving
conjugacy of the lifts.

For the state-trajectory factor, let $C_m$ be the quotient of
$\Sigma_m\times[0,1]$ obtained by collapsing $\Sigma_m\times\{0\}$ to one
point. The trajectory corresponding to $(u,s)$ is
\[
 \bigl(R^k(s),R^k(s)Y_q(\sigma^ku)\bigr)_{k\geq0}.
\]
At positive radius this trajectory determines every input label: its
successive normalized states determine $v_{u_k}(q)$, and the $v_j(q)$ are
distinct. At zero radius all inputs give the same constant trajectory.
These are exactly the identifications in $C_m$.
The trajectory map therefore induces a homeomorphism $C_m\to\T_Q$.
Under it, $S$ is the quotient of $\sigma\times R$.
The map $\mathrm{id}\times h$ passes to the quotient, proving both the
factor conjugacies and their commutation with $e$.

All feasible states at time $n$ lie in $R^n(1)Q$, whose diameter is
$2R^n(1)\to0$. The state-trajectory shift consequently converges
uniformly to its constant apex trajectory and has zero entropy: at any
fixed separation scale only finitely many initial iterates need to be
covered. In an input fiber, the normalized trajectory is fixed and only
the radial coordinate varies. Equicontinuity of $R^n$ bounds separated
sets in every fiber uniformly in time and input, proving zero fiber
entropy.
In the product model $\sigma\times R$, equicontinuity supplies a finite
radial cover valid for every iterate. Matching the first $n+\ell$ input
symbols then bounds separated sets by $Cm^{n+\ell}$ at a fixed scale,
with $C$ and $\ell$ independent of $n$. Uniform continuity of $H_{q,R}$
transfers this bound to the lift. The apex subsystem is the full shift,
so $\h(F)=\log m$.

For an invariant measure, the nonnegative continuous function
$s-R(s)$ has integral zero. It vanishes only at $s=0$ on $Q$, so the
measure is supported at the apex. Conversely every shift-invariant
$\nu$ gives $\nu\times\delta_{(0,0)}$.
On this support $\pi$ is an isomorphism, giving conditional entropy zero.
\end{proof}

\begin{corollary}\label{cor:no-go}
No functional invariant under isomorphism of the diagram
$(\Lambda_Q,F,\pi,e,\T_Q,S)$ can recover any one of the three control
entropies, already for $K=Q$, on all smooth finite-alphabet systems with compact convex
controlled invariant constraints. Adding the invariant measures of the
lift does not change this conclusion.
\end{corollary}
\begin{proof}
For $m=2$, take $(a,b)=(\log2,\log2)$ and
$(a,b)=(\log(3/2),\log(3/2))$ in \cref{thm:realization}.
The diagrams are isomorphic and their families of invariant measures
coincide, while each of the three control entropy values differs.
\end{proof}

The initial-state projection is additional data. Write
$p_Q(u,x)=x$. For a length-$n$ input word $w$, let
\[
 V_w=p_Q\bigl(\Lambda_Q\cap\pi^{-1}([w])\bigr).
\]
Controlled invariance implies that $V_w$ is exactly the set of states
whose finite trajectory under $w$ stays in $Q$: any such finite trajectory
can be continued. Therefore $\rinv(n,K,Q)$ is the minimum number of sets
$V_w$ covering $K$. The conjugacies above need not preserve these
initial-state identifications when $q$ changes.

There is a further distinction for ambient tolerances. With $q$ fixed,
the radial conjugacy induces a state homeomorphism of the constraint,
\[
 (s,sy)\longmapsto(h(s),h(s)y).
\]
It preserves every feasible controlled transition, and hence every exact
finite-horizon count on corresponding initial sets. Taking different
$b$ at fixed $a$ nevertheless changes the outer and tracking rates.
This homeomorphism is defined on $Q$; an ambient conjugacy of the
controlled maps near $Q$ has not been asserted. Thus there is no conflict
with invariance of control entropy under an ambient state equivalence.

\section{Common inputs on a control set}

\begin{proposition}\label{prop:control-set}
Consider $f(x,v)=2x+v$ on $\mathbb R$, with $U=[-1,1]$.
The interval $D=(-1,1)$ is a control set, and $Q=\overline D=[-1,1]$
is compact and controlled invariant. Each interior point is locally
accessible in both directions in one step; exact steering inside $D$
is possible, uniformly in a common finite time on each compact
subinterval of $D$. Every point of $D$ has a stationary feasible input.
For every positive-length compact $K\subset Q$,
\[
 \hinv(K,Q)=\hout(K,Q)=\htrack(K,Q)=\log2.
\]
The entire feasible lift is surjective and its input projection is a
homeomorphism to the full shift on $[-1,1]^{\mathbb N_0}$.
Nevertheless no two distinct initial states are a common-input
Li--Yorke pair, even without a feasibility requirement.
\end{proposition}
\begin{proof}
For $x,y\in D$, choose $n$ large enough that
$c=(2^nx-y)/(2^n-1)$ belongs to $D$ and apply the constant input $v=-c$.
At intermediate times
\[
 x_k=c+2^k(x-c)
     =x+\frac{2^k-1}{2^n-1}(y-x),\qquad 0\leq k\leq n.
\]
The states lie between $x$ and $y$, and $x_n=y$.
For $|x|,|y|\leq r<1$, it suffices to choose
$r(2^n+1)/(2^n-1)<1$, giving a common steering time.
Taking $v=-x$ keeps $x$ fixed forever.
The one-step forward and backward reachable intervals are
$[2x-1,2x+1]$ and $[(x-1)/2,(x+1)/2]$; each contains a neighborhood of
$x$ when $|x|<1$.

To verify maximality of $D$, a state $x>1$ satisfies
$x_k\geq1+2^k(x-1)$ under every input. At $x=1$, every forward state is
at least $1$. Such states cannot approximate any interior point of $D$.
The symmetric argument applies on the left, proving maximality among
controlled invariant approximately controllable sets.

The two inputs $v=\pm1$ give $2^n$ viable intervals of length $2^{1-n}$
that cover $Q$. For any fixed input, the final state has slope $2^n$
as a function of the initial state. Thus
\[
 |K|2^{n-1}\leq\rinv(n,K,Q)\leq2^n,
 \qquad
 \rinv(n,K,Q^\varepsilon)\geq\frac{|K|2^{n-1}}{1+\varepsilon}.
\]
These estimates give the exact and outer rates.
For tracking, choose an initial grid with mesh smaller than
$\varepsilon2^{-n}$ and use the stationary reference input $v=-c$ at
each center $c$. The errors are $2^k|x-c|$, giving an upper rate $\log2$.
The outer lower bound supplies the matching tracking rate.

For each input $u$, its unique feasible initial state is
$Y(u)=-\sum_{k\geq0}u_k/2^{k+1}$.
The proof is the same bounded-state calculation used above.
Consequently $u\mapsto(u,Y(u))$ conjugates the full input shift with the
entire lift. It is surjective because the full input shift is.
Every input fiber is a singleton, so all conditional state entropies
over input vanish.
Finally, for every common input,
\[
 |\varphi(n,x,u)-\varphi(n,y,u)|=2^n|x-y|.
\]
This precludes proximality whenever $x\ne y$.
\end{proof}

This example excludes a common-input implication based only on positive
control entropy, controllability, accessibility, and existence of recurrent
feasible controls. It does not impose recurrence on every feasible input,
and it does not exclude scrambled trajectories under different inputs.
It also excludes a conclusion required to hold for every feasible
controller: the continuous feedback $\kappa(x)=-x$ induces the identity.

For comparison, existence of a chaotic feedback can occur at zero control
information cost. Take $f(x,v)=x/2+v$, $U=[-1,1]$, and $Q=[-2,2]$.
Every input preserves $Q$, and one input spans it for every horizon, so
$\hinv(Q,Q)=\hout(Q,Q)=\htrack(Q,Q)=0$.
The last equality follows by using a finite initial grid and the common
zero input, under which initial errors decrease.
Every fixed-input map is a smooth diffeomorphism. Also $Q$ is a control
set: the constant input $v=y/2$ makes a trajectory from any $x\in Q$
converge to any prescribed $y\in Q$, while no point outside $Q$ can be
approximated from an initial state in $Q$.

Put $c=2/3$ and define $G:Q\to Q$ by $G(x)=x$ outside $[-c,c]$;
on that interval take
the continuous piecewise linear map through
\[
 (-c,-c),\quad(-c/3,c),\quad(c/3,-c),\quad(c,c).
\]
The feedback $\kappa(x)=G(x)-x/2$ is continuous and belongs to $[-1,1]$:
inside the small interval its absolute value is at most $3c/2=1$, and
outside it, within $Q$, it equals $x/2$.
The outer two inverse branches on $[-c,c]$ are
$(y-2c)/3$ and $(y+2c)/3$. Their images are disjoint, and their invariant
Cantor set is conjugate to the full two-symbol shift. Hence the induced
closed loop has an uncountable Li--Yorke scrambled set by the classical
positive-entropy theorem \cite{BGKM}.
This fixed continuous feedback gives one assignment for the entire set;
there is no pair-dependent choice. The construction does not refute an
existential implication from positive invariance entropy, but shows that
such existence alone need not reflect a positive control information cost.

\section{Relation to the existing entropy formulas}

The missing initial-state and ambient data in \cref{cor:no-go} are
essential. For hyperbolic control sets, Da Silva and Kawan use an unstable
Jacobian cocycle in their entropy formula, even when the control flow is
conjugate to the input shift \cite{DSK}. Kawan's current discrete-time
preprint likewise combines unstable expansion with conditional state
entropy \cite{Kawan}. Those formulas retain geometric information that
the conjugacy diagram in \cref{cor:no-go} omits.
The calculations here neither disprove these formulas nor supply a
variational principle for general forward constraints.
Their role in the P6 project is to identify what a proposed general
principle or state-chaos implication must retain.

\begin{thebibliography}{9}
\bibitem{BGKM}
F. Blanchard, E. Glasner, S. Kolyada, and A. Maass,
On Li-Yorke pairs, \emph{Journal f\"ur die reine und angewandte Mathematik}
\textbf{547} (2002), 51--68.
\href{https://doi.org/10.1515/crll.2002.053}{doi:10.1515/crll.2002.053}.
\bibitem{DSK}
A. Da Silva and C. Kawan, Invariance entropy of hyperbolic control sets,
\emph{Discrete and Continuous Dynamical Systems} \textbf{36} (2016),
no.~1, 97--136.
\href{https://doi.org/10.3934/dcds.2016.36.97}{doi:10.3934/dcds.2016.36.97}.
\bibitem{Kawan}
C. Kawan, \emph{Control of chaos with minimal information transfer},
arXiv:2003.06935, version 5, 20 September 2026.
\href{https://arxiv.org/abs/2003.06935v5}{Version 5};
\href{https://doi.org/10.48550/arXiv.2003.06935}{arXiv DOI}.
\end{thebibliography}
\end{document}
